% Présentation 3 : Philosophie d'enseignement et offre de cours
% Carlos Denner dos Santos
% Journée d'évaluation - 23 janvier 2026
% Poste de professeur en gestion de l'intelligence artificielle
% Département SIMQG, École de gestion, Université de Sherbrooke

% === PDF SETTINGS (full font embedding) ===
\pdfminorversion=7
\pdfobjcompresslevel=0

\documentclass[aspectratio=169,11pt]{beamer}

% === ENCODAGE ET LANGUE ===
\usepackage[utf8]{inputenc}
\usepackage[T1]{fontenc}
\usepackage[french]{babel}
\usepackage{csquotes}
\usepackage{lmodern}  % Latin Modern fonts (fully embeddable)

% === STYLE UdeS ===
\usepackage{beamer_usherbrooke}

% === PACKAGES ADDITIONNELS ===
\usepackage{graphicx}
\usepackage{booktabs}
\usepackage{array}
\usepackage{colortbl}
\usepackage{tikz}
\usetikzlibrary{shapes,arrows,positioning,calc}
\usepackage{hyperref}

% === COULEURS PERSONNALISÉES ===
\definecolor{immediatgreen}{RGB}{34,139,34}   % Vert pour (I)
\definecolor{prepareorange}{RGB}{255,140,0}   % Orange pour (P)

% === MÉTADONNÉES ===
\title[Philosophie d'enseignement]{Philosophie d'enseignement et offre de cours}
\subtitle{Poste de professeur en gestion de l'intelligence artificielle}
\author[C. D. dos Santos]{Carlos Denner dos Santos Jr.}
\institute[UdeS - SIMQG]{
  Département des systèmes d'information\\
  et méthodes quantitatives de gestion (SIMQG)\\
  École de gestion, Université de Sherbrooke
}
\date{23 janvier 2026}

\begin{document}

% ============================================
% SLIDE TITRE
% ============================================
\begin{frame}
  \begin{beamercolorbox}[wd=\paperwidth,ht=2.5cm,dp=0.5cm,center]{palette primary}
    {\Large\bfseries Présentation 3 : Philosophie d'enseignement}
  \end{beamercolorbox}
  
  \vspace{0.8cm}
  
  \begin{center}
    {\Large Carlos Denner dos Santos Jr.}
    
    \vspace{0.8em}
    
    {\normalsize Département des systèmes d'information\\
    et méthodes quantitatives de gestion (SIMQG)\\
    École de gestion, Université de Sherbrooke}
    
    \vspace{0.8em}
    
    {\normalsize 23 janvier 2026}
    
    \vspace{1em}
    
    {\large\textcolor{ushervert}{\textbf{\guillemotleft~Je forme des gestionnaires qui savent décider avec l'IA.~\guillemotright}}}
  \end{center}
\end{frame}

% ============================================
% PLAN DE LA PRÉSENTATION
% ============================================
\begin{frame}{Plan de la présentation}
  \begin{columns}[T]
    \begin{column}{0.48\textwidth}
      \begin{block}{A. Philosophie pédagogique}
        \begin{itemize}
          \item Une anecdote formatrice
          \item Quatre principes directeurs
          \item Alignement avec l'UdeS
          \item Approche d'évaluation
        \end{itemize}
      \end{block}
    \end{column}
    \begin{column}{0.48\textwidth}
      \begin{block}{B. Offre de cours}
        \begin{itemize}
          \item B.A.A. -- Cheminement GTA
          \item M. Sc. -- Commerce électronique
          \item M. Sc. -- Stratégie IA (SIA)
          \item Microprogramme 3\textsuperscript{e} cycle IA
        \end{itemize}
      \end{block}
    \end{column}
  \end{columns}
\end{frame}

% ============================================
% PARTIE A : PHILOSOPHIE PÉDAGOGIQUE
% ============================================
\begin{frame}{}
  \begin{beamercolorbox}[wd=\paperwidth,ht=2.5cm,dp=0.5cm,center]{palette primary}
    {\Large\bfseries Partie A : Philosophie pédagogique}
  \end{beamercolorbox}
  \vspace{2cm}
  \begin{center}
    {\large\textit{Former des gestionnaires capables de penser,\\d'agir et d'innover avec l'IA}}
  \end{center}
\end{frame}

% --- SLIDE : Anecdote formatrice ---
\begin{frame}{Une anecdote formatrice}
  \begin{columns}[T]
    \begin{column}{0.55\textwidth}
      \begin{block}{Le projet qui a tout changé}
        \textbf{Brasilia, 2018} -- Cours de BI avancée
        
        \vspace{0.5em}
        Un groupe d'étudiants décroche un partenariat avec l'\emphvert{hôpital universitaire} pour analyser ses données patients.
        
        \vspace{0.5em}
        Problème : les données sont \emphor{incomplètes, biaisées}, et les résultats de leur modèle posent des \emphor{questions éthiques}.
      \end{block}
    \end{column}
    \begin{column}{0.42\textwidth}
      \begin{alertblock}{La leçon}
        Ce projet m'a convaincu que :
        \begin{itemize}
          \item Les \emphvert{cas réels} sont irremplaçables
          \item La \emphvert{technique seule} ne suffit pas
          \item L'\emphvert{éthique} doit être intégrée dès le départ
        \end{itemize}
      \end{alertblock}
    \end{column}
  \end{columns}
  
  \vspace{0.5em}
  \begin{center}
    \textit{``Depuis, chaque cours commence par une situation authentique.''}
  \end{center}
\end{frame}

% --- SLIDE : 4 Principes pédagogiques (Vue d'ensemble) ---
\begin{frame}{Quatre principes directeurs}
  \begin{center}
    \begin{tikzpicture}[scale=0.9]
      % Principe 1 - Situations authentiques
      \node[draw, fill=ushervert!20, rounded corners, minimum width=5cm, minimum height=1.2cm, align=center] (p1) at (-4,2) {\textbf{1. Situations authentiques}\\{\small Partir de problèmes réels}};
      
      % Principe 2 - Learning by doing
      \node[draw, fill=ushervert!20, rounded corners, minimum width=5cm, minimum height=1.2cm, align=center] (p2) at (4,2) {\textbf{2. Learning by doing}\\{\small Équilibre théorie-pratique}};
      
      % Principe 3 - Autonomie réflexive
      \node[draw, fill=ushervert!20, rounded corners, minimum width=5cm, minimum height=1.2cm, align=center] (p3) at (-4,-1) {\textbf{3. Autonomie réflexive}\\{\small Penser par soi-même}};
      
      % Principe 4 - Climat inclusif
      \node[draw, fill=ushervert!20, rounded corners, minimum width=5cm, minimum height=1.2cm, align=center] (p4) at (4,-1) {\textbf{4. Climat inclusif}\\{\small Bienveillant mais rigoureux}};
      
      % Centre
      \node[draw, fill=usheror!30, circle, minimum size=2.5cm, align=center, font=\small] (centre) at (0,0.5) {\textbf{Étudiant}\\\textbf{autonome}};
      
      % Flèches
      \draw[->, thick, ushervert] (p1) -- (centre);
      \draw[->, thick, ushervert] (p2) -- (centre);
      \draw[->, thick, ushervert] (p3) -- (centre);
      \draw[->, thick, ushervert] (p4) -- (centre);
    \end{tikzpicture}
  \end{center}
\end{frame}

% --- SLIDE : Principe 1 - Situations authentiques ---
\begin{frame}{Principe 1 : Ancrer l'apprentissage dans le réel}
  \begin{columns}[T]
    \begin{column}{0.55\textwidth}
      \begin{block}{Approche}
        \begin{itemize}
          \item Commencer par un \emphvert{problème concret}
          \item Exemple : un projet IA mal cadré, une startup qui échoue, un biais algorithmique
          \item Les concepts viennent \textbf{après} -- pour résoudre le problème
        \end{itemize}
      \end{block}
      
      \vspace{0.3em}
      
      \begin{exampleblock}{En pratique}
        \small
        \textit{``Voici les données d'une entreprise qui a investi 2M\$ en IA sans retour. Que s'est-il passé?''}
      \end{exampleblock}
    \end{column}
    \begin{column}{0.42\textwidth}
      \begin{alertblock}{Pourquoi ça fonctionne}
        \begin{itemize}
          \item \emphvert{Motivation} intrinsèque
          \item Ancrage mémoriel fort
          \item Transfert facilité vers le milieu professionnel
        \end{itemize}
      \end{alertblock}
      
      \vspace{0.3em}
      
      {\small Aligné avec la \emphvert{pédagogie active} de l'École de gestion}
    \end{column}
  \end{columns}
\end{frame}

% --- SLIDE : Principe 2 - Learning by doing ---
\begin{frame}{Principe 2 : Combiner rigueur conceptuelle et pratique}
  \begin{columns}[T]
    \begin{column}{0.48\textwidth}
      \begin{block}{L'équilibre théorie-pratique}
        \begin{itemize}
          \item Chaque concept $\rightarrow$ exercice concret
          \item Études de cas, mini-projets, prototypage
          \item Itérations : faire, échouer, apprendre
        \end{itemize}
      \end{block}
      
      \vspace{0.3em}
      
      \begin{exampleblock}{Exemples d'activités}
        \small
        \begin{itemize}
          \item Notebook Python en labo
          \item Audit d'un système IA existant
          \item Pitch d'un projet IA (5 min)
        \end{itemize}
      \end{exampleblock}
    \end{column}
    \begin{column}{0.48\textwidth}
      \begin{alertblock}{Innovation pédagogique}
        J'ai publié sur l'innovation pédagogique :
        \begin{itemize}
          \item Exercices en classe interactifs
          \item Projets itératifs avec feedback
          \item \emphvert{Jeux sérieux} et simulations
        \end{itemize}
        
        \vspace{0.2em}
        
        {\scriptsize \textit{Apprentissage par l'expérience et la ludification}}
      \end{alertblock}
      
      \vspace{0.3em}
      
      {\small Les étudiants apprennent les \emphvert{limites} des outils en les utilisant}
    \end{column}
  \end{columns}
\end{frame}

% --- SLIDE : Principe 3 - Autonomie réflexive ---
\begin{frame}{Principe 3 : Cultiver l'autonomie réflexive}
  \begin{columns}[T]
    \begin{column}{0.55\textwidth}
      \begin{block}{Former des penseurs critiques}
        \begin{itemize}
          \item Poser des \emphvert{questions ouvertes}
          \item Pourquoi ce modèle? Quels biais? Quelles implications éthiques?
          \item Développer le jugement, pas seulement la technique
        \end{itemize}
      \end{block}
      
      \vspace{0.3em}
      
      \begin{exampleblock}{Question type en classe}
        \small
        \textit{``Votre modèle prédit un risque de défaut de 78\%. Feriez-vous confiance à cette prédiction si le client était votre mère?''}
      \end{exampleblock}
    \end{column}
    \begin{column}{0.42\textwidth}
      \begin{alertblock}{Éthique intégrée}
        J'intègre \textbf{systématiquement} :
        \begin{itemize}
          \item Questions de \emphvert{biais}
          \item Enjeux de \emphvert{transparence}
          \item Risques de \emphvert{sécurité}
          \item Implications \emphvert{sociétales}
        \end{itemize}
      \end{alertblock}
      
      \vspace{0.3em}
      
      {\small Former des \emphvert{gestionnaires responsables}}
    \end{column}
  \end{columns}
\end{frame}

% --- SLIDE : Principe 4 - Climat inclusif ---
\begin{frame}{Principe 4 : Construire un climat inclusif et exigeant}
  \begin{columns}[T]
    \begin{column}{0.48\textwidth}
      \begin{block}{Bienveillant mais rigoureux}
        \begin{itemize}
          \item Attentes \emphvert{claires} dès le départ
          \item Feedback \emphvert{régulier} et constructif
          \item Droit à l'erreur (c'est ainsi qu'on apprend)
        \end{itemize}
      \end{block}
      
      \vspace{0.3em}
      
      \begin{exampleblock}{Adaptation aux publics}
        \small
        \begin{itemize}
          \item Étudiants en emploi
          \item Étudiants internationaux
          \item Parcours atypiques
        \end{itemize}
      \end{exampleblock}
    \end{column}
    \begin{column}{0.48\textwidth}
      \begin{alertblock}{Mon expérience}
        \begin{itemize}
          \item \textbf{15+ ans} d'enseignement
          \item Brésil et États-Unis
          \item Présentiel et en ligne
          \item Tous les cycles (bacc. à doctorat)
        \end{itemize}
      \end{alertblock}
      
      \vspace{0.3em}
      
      {\small Cette diversité m'a rendu \emphvert{sensible} à l'adaptation pédagogique}
    \end{column}
  \end{columns}
\end{frame}

% --- SLIDE : Alignement UdeS ---
\begin{frame}{Alignement avec les valeurs de l'UdeS}
  \begin{center}
    \renewcommand{\arraystretch}{1.4}
    \begin{tabular}{>{\raggedright}p{4.5cm}|>{\raggedright\arraybackslash}p{7cm}}
      \rowcolor{ushervert}
      \textcolor{white}{\textbf{Valeur UdeS}} & \textcolor{white}{\textbf{Ma pratique}} \\
      \hline
      \rowcolor{ushervertclair}
      \textbf{Pédagogie active} & Cas réels, projets, prototypage \\
      \hline
      \textbf{Engagement étudiant} & Questions ouvertes, participation active \\
      \hline
      \rowcolor{ushervertclair}
      \textbf{Interdisciplinarité} & Gestion + Informatique + Statistique + Éthique \\
      \hline
      \textbf{Innovation} & LLM en classe, ludification, simulations \\
      \hline
      \rowcolor{ushervertclair}
      \textbf{Responsabilité} & Éthique et gouvernance intégrées \\
    \end{tabular}
  \end{center}
  
  \vspace{0.5em}
  
  \begin{center}
    \emphvert{Ma philosophie est naturellement alignée avec la mission de l'École de gestion}
  \end{center}
\end{frame}

% --- SLIDE : Approche d'évaluation ---
\begin{frame}{Approche d'évaluation centrée sur l'apprentissage}
  \begin{columns}[T]
    \begin{column}{0.48\textwidth}
      \begin{block}{Diversité des évaluations}
        \begin{itemize}
          \item Travaux d'\emphvert{équipe} (collaboration)
          \item Travaux \emphvert{individuels} (autonomie)
          \item \emphvert{Coévaluation} par les pairs
          \item Présentations orales
        \end{itemize}
      \end{block}
      
      \vspace{0.3em}
      
      \begin{exampleblock}{Feedback constructif}
        \small
        \begin{itemize}
          \item Commentaires détaillés
          \item Points forts et axes d'amélioration
          \item Opportunités de révision
        \end{itemize}
      \end{exampleblock}
    \end{column}
    \begin{column}{0.48\textwidth}
      \begin{alertblock}{Objectif}
        \textbf{L'évaluation au service de l'apprentissage}
        
        \vspace{0.3em}
        
        \begin{itemize}
          \item Mesurer les \emphvert{compétences} (pas seulement la mémorisation)
          \item Préparer au \emphvert{monde professionnel}
          \item Favoriser la \emphvert{réussite étudiante}
        \end{itemize}
      \end{alertblock}
      
      \vspace{0.2em}
      
      \begin{block}{Équité et intégrité}
        \scriptsize
        \begin{itemize}
          \item Grilles + rubriques + calibration
          \item IA générative : règles + traçabilité\\{\tiny (je note la démarche, pas le texte final)}
        \end{itemize}
      \end{block}
    \end{column}
  \end{columns}
\end{frame}

% ============================================
% PARTIE B : OFFRE DE COURS
% ============================================
\begin{frame}{}
  \begin{beamercolorbox}[wd=\paperwidth,ht=2.5cm,dp=0.5cm,center]{palette primary}
    {\Large\bfseries Partie B : Offre de cours}
  \end{beamercolorbox}
  \vspace{1.5cm}
  \begin{center}
    {\large\textit{Une contribution immédiate aux programmes de l'École}}
    
    \vspace{1em}
    
    \begin{tabular}{cl}
      \textcolor{immediatgreen}{\textbf{(I)}} & Immédiat -- prêt à enseigner \\
      \textcolor{prepareorange}{\textbf{(P)}} & Avec préparation -- mise à jour ciblée \\
    \end{tabular}
  \end{center}
\end{frame}

% --- SLIDE : B.A.A. GTA ---
\begin{frame}{B.A.A. -- Cheminement Gestion des technologies d'affaires}
  \small
  \begin{columns}[T]
    \begin{column}{0.48\textwidth}
      \begin{block}{Cours en (I) Immédiat}
        \begin{itemize}
          \item[\textcolor{immediatgreen}{I}] \textbf{GTA211} Technologies d'affaires
          \item[\textcolor{immediatgreen}{I}] \textbf{GTA311} Valorisation des données
          \item[\textcolor{immediatgreen}{I}] \textbf{GTA431} Manipulation des données
          \item[\textcolor{immediatgreen}{I}] \textbf{GTA651} \emphvert{IA appliquée}
        \end{itemize}
      \end{block}
    \end{column}
    \begin{column}{0.48\textwidth}
      \begin{block}{Cours en (P) Préparation}
        \begin{itemize}
          \item[\textcolor{prepareorange}{P}] \textbf{GTA411} Gestion de projets techno.
          \item[\textcolor{prepareorange}{P}] \textbf{GTA531} Transformation numérique
        \end{itemize}
      \end{block}
      
      \vspace{0.5em}
      
      \begin{alertblock}{Bilan GTA}
        \centering
        \textbf{4 cours (I)} sur 6 cours ciblés
        
        \vspace{0.3em}
        
        {\small Couverture des fondamentaux du cheminement}
      \end{alertblock}
    \end{column}
  \end{columns}
\end{frame}

% --- SLIDE : M. Sc. Commerce électronique ---
\begin{frame}{M. Sc. -- Gestion du commerce électronique}
  \small
  \begin{columns}[T]
    \begin{column}{0.48\textwidth}
      \begin{block}{Cours en (I) Immédiat}
        \begin{itemize}
          \item[\textcolor{immediatgreen}{I}] \textbf{GCE806} Stratégie numérique
          \item[\textcolor{immediatgreen}{I}] \textbf{GCE821} Planif. projets numériques
          \item[\textcolor{immediatgreen}{I}] \textbf{GCE846} Web analytique
          \item[\textcolor{immediatgreen}{I}] \textbf{GTA821} Technologies entreprise num.
        \end{itemize}
      \end{block}
    \end{column}
    \begin{column}{0.48\textwidth}
      \begin{block}{Cours en (P) Préparation}
        \begin{itemize}
          \item[\textcolor{prepareorange}{P}] \textbf{GCE802} Fondements e-commerce
          \item[\textcolor{prepareorange}{P}] \textbf{GCE811} Enjeux légaux du numérique
          \item[\textcolor{prepareorange}{P}] \textbf{GCE826} Création sites Web
        \end{itemize}
      \end{block}
      
      \vspace{0.3em}
      
      {\small \textit{Volets marketing/UX : co-enseignement en complémentarité}}
    \end{column}
  \end{columns}
  
  \vspace{0.3em}
  
  \begin{center}
    \emphvert{Couverture complète des volets technologiques et analytiques}
  \end{center}
\end{frame}

% --- SLIDE : M. Sc. SIA ---
\begin{frame}{M. Sc. -- Stratégie de l'intelligence d'affaires (SIA)}
  \scriptsize
  \begin{columns}[T]
    \begin{column}{0.48\textwidth}
      \begin{block}{Cours en (I) Immédiat -- 7 cours}
        \begin{itemize}
          \item[\textcolor{immediatgreen}{I}] \textbf{GIS800} Fondements de l'IA
          \item[\textcolor{immediatgreen}{I}] \textbf{GIS801} IA stratégique
          \item[\textcolor{immediatgreen}{I}] \textbf{GIS807} Tableaux de bord et BI
          \item[\textcolor{immediatgreen}{I}] \textbf{GIS811} Architectures données IA
          \item[\textcolor{immediatgreen}{I}] \textbf{GIS812} Analytique et science données
          \item[\textcolor{immediatgreen}{I}] \textbf{GIS861} Méthodes de recherche IA
          \item[\textcolor{immediatgreen}{I}] \textbf{MQG811} Analyse statistique
        \end{itemize}
      \end{block}
    \end{column}
    \begin{column}{0.48\textwidth}
      \begin{block}{Cours en (P) Préparation -- 3 cours}
        \begin{itemize}
          \item[\textcolor{prepareorange}{P}] \textbf{GIS806} Entrepôts de données
          \item[\textcolor{prepareorange}{P}] \textbf{MQG812} Forage de données
          \item[\textcolor{prepareorange}{P}] \textbf{MQG813} Stats décisionnelles avancées
        \end{itemize}
      \end{block}
      
      \vspace{0.5em}
      
      \begin{alertblock}{Bilan SIA}
        \centering
        \textbf{7 cours (I)} sur 12 cours
        
        \vspace{0.3em}
        
        {\small Programme de \emphvert{prédilection}\\(analytique + IA + gestion)}
      \end{alertblock}
    \end{column}
  \end{columns}
  
  \vspace{0.3em}
  
  \begin{center}
    {\small Mon postdoc en \emphvert{statistique} + expérience en \emphvert{data science} = couverture quasi-complète}
  \end{center}
\end{frame}

% --- SLIDE : Microprogramme 3e cycle IA ---
\begin{frame}{Microprogramme de 3\textsuperscript{e} cycle -- Gestion stratégique de l'IA}
  \begin{columns}[T]
    \begin{column}{0.55\textwidth}
      \begin{block}{5 cours obligatoires -- Tous en (I)}
        \begin{itemize}
          \item[\textcolor{immediatgreen}{I}] \textbf{GTA901} Preuve de concept en IA
          \item[\textcolor{immediatgreen}{I}] \textbf{GTA902} Valeur d'affaires de l'IA
          \item[\textcolor{immediatgreen}{I}] \textbf{GTA903} Faisabilité de l'IA
          \item[\textcolor{immediatgreen}{I}] \textbf{GTA904} Planification implantation IA
          \item[\textcolor{immediatgreen}{I}] \textbf{GTA905} \emphvert{Gouvernance de l'IA}
        \end{itemize}
      \end{block}
      
      \vspace{0.3em}
      
      {\small Ce programme est \emphvert{directement aligné} avec mon expertise et mes recherches}
    \end{column}
    \begin{column}{0.42\textwidth}
      \begin{alertblock}{Valeur ajoutée}
        \begin{itemize}
          \item Cas récents (sécurité LLM, Firewall)
          \item Frameworks de gouvernance publiés
          \item Expérience terrain (Bell, Videns)
        \end{itemize}
      \end{alertblock}
      
      \vspace{0.5em}
      
      \begin{exampleblock}{Bilan}
        \centering
        \textbf{5/5 cours en (I)}
        
        \vspace{0.3em}
        
        {\small Pleinement autonome}
      \end{exampleblock}
    \end{column}
  \end{columns}
\end{frame}

% --- SLIDE : Synthèse des cours ---
\begin{frame}{Synthèse : Capacité d'enseignement}
  \begin{center}
    \renewcommand{\arraystretch}{1.3}
    \begin{tabular}{l|c|c|c}
      \rowcolor{ushervert}
      \textcolor{white}{\textbf{Programme}} & \textcolor{white}{\textbf{(I) Immédiat}} & \textcolor{white}{\textbf{(P) Prép.}} & \textcolor{white}{\textbf{Total}} \\
      \hline
      \rowcolor{ushervertclair}
      B.A.A. -- GTA & \textcolor{immediatgreen}{\textbf{4}} & \textcolor{prepareorange}{2} & 6 \\
      \hline
      M. Sc. -- Commerce élec. & \textcolor{immediatgreen}{\textbf{4}} & \textcolor{prepareorange}{3} & 7 \\
      \hline
      \rowcolor{ushervertclair}
      M. Sc. -- SIA & \textcolor{immediatgreen}{\textbf{7}} & \textcolor{prepareorange}{3} & 10 \\
      \hline
      Microprog. 3\textsuperscript{e} cycle IA & \textcolor{immediatgreen}{\textbf{5}} & \textcolor{prepareorange}{0} & 5 \\
      \hline
      \rowcolor{usheror!20}
      \textbf{TOTAL} & \textcolor{immediatgreen}{\textbf{20}} & \textcolor{prepareorange}{8} & \textbf{28} \\
    \end{tabular}
  \end{center}
  
  \vspace{0.3em}
  
  \begin{center}
    \emphvert{20 cours prêts immédiatement} -- mais mon objectif : \textbf{2--3 cours signature}\\[0.3em]
    {\small Gouvernance opérationnelle de l'IA $\cdot$ Sécurité LLM $\cdot$ Éthique appliquée}
  \end{center}
\end{frame}

% --- SLIDE : Au-delà des cours existants ---
\begin{frame}{Au-delà des cours existants}
  \begin{columns}[T]
    \begin{column}{0.48\textwidth}
      \begin{block}{Créer et innover}
        \begin{itemize}
          \item Co-développer de \emphvert{nouveaux cours}
          \item Actualiser les cours existants avec mes \emphvert{études de cas}
          \item Intégrer les \emphvert{dernières recherches} (LLM, gouvernance)
        \end{itemize}
      \end{block}
      
      \vspace{0.3em}
      
      \begin{exampleblock}{Idées de nouveaux cours}
        \small
        \begin{itemize}
          \item \textbf{Gouvernance opérationnelle de l'IA}\\
                {\scriptsize (cadres + contrôles + cas réels)}
          \item Sécurité des systèmes LLM
          \item Éthique appliquée de l'IA
        \end{itemize}
      \end{exampleblock}
    \end{column}
    \begin{column}{0.48\textwidth}
      \begin{alertblock}{Encadrement}
        \begin{itemize}
          \item Travaux dirigés
          \item Mémoires et thèses
          \item Projets d'intégration
        \end{itemize}
        
        \vspace{0.3em}
        
        {\small Déjà $\sim$\textbf{9 doctorats} et $\sim$\textbf{6 maîtrises} dirigés}
      \end{alertblock}
      
      \vspace{0.3em}
      
      \begin{block}{Contribuer à}
        \small
        \begin{itemize}
          \item Firme-école
          \item Compétitions étudiantes
          \item École d'été en IA stratégique
        \end{itemize}
      \end{block}
    \end{column}
  \end{columns}
\end{frame}

% ============================================
% CONCLUSION
% ============================================
\begin{frame}{}
  \begin{beamercolorbox}[wd=\paperwidth,ht=2.5cm,dp=0.5cm,center]{palette primary}
    {\Large\bfseries Conclusion}
  \end{beamercolorbox}
  \vspace{2cm}
  \begin{center}
    {\large\textit{Former des leaders du numérique\\alliant technique, stratégie et éthique}}
  \end{center}
\end{frame}

% --- SLIDE : En résumé ---
\begin{frame}{En résumé : Pourquoi moi pour l'enseignement?}
  \begin{columns}[T]
    \begin{column}{0.55\textwidth}
      \begin{enumerate}
        \item[1.] \emphvert{Philosophie alignée UdeS}\\
        {\small Pédagogie active, cas réels, éthique intégrée}
        \vspace{0.25em}
        \item[2.] \emphvert{Expérience démontrée}\\
        {\small 15+ ans, tous cycles, 2 pays, présentiel/ligne}
        \vspace{0.25em}
        \item[3.] \emphvert{Couverture large}\\
        {\small 20 cours (I), 4 programmes, $\sim$15 étudiants gradués}
        \vspace{0.25em}
        \item[4.] \emphvert{Innovation pédagogique}\\
        {\small LLM en classe, projets itératifs, coévaluation}
        \vspace{0.25em}
        \item[5.] \emphvert{Engagement total}\\
        {\small Encadrement, firme-école, École d'été IA}
      \end{enumerate}
    \end{column}
    \begin{column}{0.4\textwidth}
      \centering
      \colorbox{ushervert}{\parbox{3.8cm}{\centering\textcolor{white}{\large\bfseries Prêt à contribuer\\dès le jour 1}}}
      
      \vspace{0.5em}
      
      {\scriptsize\textit{12 premiers mois : plusieurs cours + un cas d'enseignement issu de mes terrains}}
      
      \vspace{0.5em}
      
      {\small\textbf{Je forme des diplômés capables de décider, pas juste d'utiliser des outils.}}
    \end{column}
  \end{columns}
\end{frame}

% --- SLIDE FINAL ---
\begin{frame}{}
  \begin{beamercolorbox}[wd=\paperwidth,ht=1.8cm,dp=0.3cm,center]{palette primary}
    {\large\bfseries Merci de votre attention}
  \end{beamercolorbox}
  
  \vspace{1cm}
  
  \centering
  {\Huge\emphvert{Merci!}}
  
  \vspace{0.8em}
  
  {\Large Questions?}
  
  \vspace{1em}
  
  \begin{tabular}{cl}
    Courriel : & carlosdenner@gmail.com \\
    Téléphone : & +1 (438) 836-4116 \\
    LinkedIn : & linkedin.com/in/carlosdenner \\
    GitHub : & github.com/carlosdenner \\
  \end{tabular}
\end{frame}

\end{document}
